
\documentclass[25pt,a0paper,landscape]{tikzposter}

\usepackage[T1]{fontenc}
\usepackage[utf8]{inputenc}
\usepackage{lmodern}
\usepackage[scaled]{helvet}
\renewcommand{\familydefault}{\sfdefault}

\usepackage{amsmath,amssymb,booktabs,array,multirow}
\usepackage{graphicx}
\usepackage[table]{xcolor}
\usepackage{enumitem}
\usepackage{tikz}
\usetikzlibrary{positioning,arrows.meta,calc}
\usepackage{ragged2e}

\usetheme{Simple}

\definecolor{AccentBlue}{HTML}{0B5FA5}
\definecolor{AccentTeal}{HTML}{0F8B8D}
\definecolor{SoftBlue}{HTML}{EAF3FF}
\definecolor{SoftTeal}{HTML}{E9F9F7}
\definecolor{SoftGray}{HTML}{F5F7FA}
\definecolor{DarkGray}{HTML}{2B2B2B}

\colorlet{backgroundcolor}{white}
\colorlet{framecolor}{AccentBlue}
\colorlet{titlefgcolor}{white}
\colorlet{titlebgcolor}{AccentBlue}
\colorlet{blocktitlebgcolor}{AccentBlue}
\colorlet{blocktitlefgcolor}{white}
\colorlet{blockbodybgcolor}{SoftGray}
\colorlet{blockbodyfgcolor}{DarkGray}
\colorlet{innerblocktitlebgcolor}{AccentTeal}
\colorlet{innerblocktitlefgcolor}{white}
\colorlet{innerblockbodybgcolor}{SoftTeal}
\colorlet{innerblockbodyfgcolor}{DarkGray}
\colorlet{notefgcolor}{DarkGray}
\colorlet{notebgcolor}{SoftBlue}
\colorlet{noteframecolor}{AccentBlue}

\setlength{\parindent}{0pt}
\setlength{\parskip}{0.35em}
\setlist[itemize]{leftmargin=1.2em,itemsep=0.25em,topsep=0.25em}

\def\TP@titletextscale{0.68}
% You can replace this with the exact paper title if you prefer.
\title{MEDICAL ASR WITH MEDICAL-AWARE REINFORCEMENT FINE-TUNING}
\author{Congjie Wang (1) \quad Xiaofan Ye (4) \quad Jinlin Wu (2,3) \quad Dong Yi (2) \quad Zhen Lei (2,3) \quad Wai S.~Poon (4) \quad Hongbin Liu (2)}
\institute{(1) The University of Hong Kong \qquad (2) CAIR, Chinese Academy of Sciences \qquad (3) MAIS, Institute of Automation, CAS \qquad (4) HKU-Shenzhen Hospital}

\begin{document}
\maketitle

\block{TAKE-HOME MESSAGE}{
\textbf{Problem:} Standard ASR under-weights rare but clinically critical medical terms. \quad
\textbf{Solution:} Medical-aware reinforcement fine-tuning with multi-hypothesis decoding, GPT-4o term proposal, UMLS grounding, and DPO. \quad
\textbf{Outcome:} Better medical term recognition with strong cross-dataset generalization.
}

\block{PROPOSED FRAMEWORK}{
\begin{itemize}
    \item \textbf{Step 1:} generate multiple candidate transcriptions for each utterance using diverse Whisper decoding strategies.
    \item \textbf{Step 2:} use GPT-4o to propose candidate medical terms from each transcript.
    \item \textbf{Step 3:} validate and normalize the extracted terms with UMLS.
    \item \textbf{Step 4:} score each hypothesis with a composite reward combining WER, MWER, and length regularization.
    \item \textbf{Step 5:} create preference pairs and update the ASR model with DPO.
\end{itemize}
}

\begin{columns}
\column{0.32}

\block{MOTIVATION \& CONTRIBUTIONS}{
\begin{itemize}
    \item \textbf{Why medical ASR is hard:} fluent transcripts can still contain dangerous errors in drug names, dosages, anatomy, and abbreviations.
    \item \textbf{Key limitation of standard training:} sequence-level cross-entropy optimizes every token equally, so rare domain-critical terms are under-emphasized.
    \item \textbf{Our contributions:}
    \begin{itemize}
        \item medical-aware reinforcement fine-tuning (RFT) for ASR;
        \item a composite reward that combines \textbf{WER + MWER + length regularization};
        \item GPT-4o high-recall term proposal with \textbf{UMLS} precision filtering;
        \item preference-based optimization that improves specialized vocabulary while maintaining fluency.
    \end{itemize}
\end{itemize}

\textbf{Design principle:} use domain knowledge only during training; at inference time, decode with a standard single-hypothesis ASR model.
}

\block{MEDICAL-AWARE REWARD}{
For a hypothesis $h_i$ and reference $r$, the reward is
\[
R(h_i,r) = (1-\lambda)\big[\mathrm{WERScore}(h_i,r)\cdot L(h_i,r)\big] + \lambda\,\mathrm{MedScore}(h_i,r),
\]
where
\[
\mathrm{WERScore}(h_i,r)=1-\min(\mathrm{WER}(h_i,r),1.0),
\]
\[
\mathrm{MedScore}(h_i,r)=1-\min(\mathrm{MWER}(h_i,r),1.0),
\]
\[
L(h_i,r)=\left(\min\left(\frac{|h_i|}{|r|},\frac{|r|}{|h_i|}\right)\right)^{0.5}.
\]

\begin{itemize}
    \item \textbf{WERScore} preserves general transcription quality.
    \item \textbf{MedScore} explicitly rewards accurate medical terminology after UMLS filtering.
    \item \textbf{Length factor} suppresses abnormal-length outputs and stabilizes preference construction.
\end{itemize}
}

\column{0.32}

\block{PREFERENCE OPTIMIZATION}{
\begin{itemize}
    \item Generate $K$ hypotheses for each utterance.
    \item Rank them using the composite reward.
    \item Build a preference pair $(h^+, h^-)$ only when the reward gap exceeds a margin $\delta = 0.1$.
    \item Optimize with \textbf{Direct Preference Optimization (DPO)} against a frozen reference policy.
\end{itemize}

\[
\mathcal{L}_{\mathrm{DPO}}
=
-\log \sigma\!\left(
\beta \log \frac{\pi(h^+|x)}{\pi_{\mathrm{ref}}(h^+|x)}
-
\beta \log \frac{\pi(h^-|x)}{\pi_{\mathrm{ref}}(h^-|x)}
\right)
\]

\begin{itemize}
    \item This directly encourages the ASR model to prefer hypotheses with better clinical content.
    \item The method avoids relying on a post-hoc LLM correction stage at inference time.
\end{itemize}
}

\block{EXPERIMENTAL SETUP}{
\textbf{Backbone:} Whisper-small

\textbf{Training data:} English split of \textbf{MultiMed}

\textbf{External evaluation:}
\begin{itemize}
    \item \textbf{Ankit} English medical ASR test split
    \item \textbf{Macabdul} MedicalASR evaluation set
\end{itemize}

\textbf{Implementation details:}
\begin{itemize}
    \item GPT-4o assists only with medical-term extraction during training
    \item DPO scale $\beta = 0.1$
    \item learning rate $= 1\times10^{-5}$
    \item $\lambda$ gradually increased from $0.1$ to $0.3$
    \item audio length capped at 30 seconds
\end{itemize}

\textbf{Metrics:}
\begin{itemize}
    \item \textbf{WER}: word error rate
    \item \textbf{MWER}: medical word error rate
\end{itemize}
}

\column{0.32}

\block{RESULTS}{
\centering
\renewcommand{\arraystretch}{1.18}
\resizebox{\linewidth}{!}{
\begin{tabular}{lcccccc}
\toprule
\multirow{2}{*}{\textbf{Model}} & \multicolumn{2}{c}{\textbf{MultiMed}} & \multicolumn{2}{c}{\textbf{Ankit}} & \multicolumn{2}{c}{\textbf{Macabdul}} \\
 & \textbf{WER} & \textbf{MWER} & \textbf{WER} & \textbf{MWER} & \textbf{WER} & \textbf{MWER} \\
\midrule
Whisper & 0.2140 & 0.1243 & 0.2018 & 0.1172 & 0.1722 & 0.0930 \\
Whisper + GPT-4o & 0.1780 & 0.1094 & 0.1990 & 0.1107 & 0.1654 & 0.0916 \\
SFT & 0.1517 & 0.0934 & 0.2140 & 0.1323 & 0.2013 & 0.1149 \\
RFT ($\lambda=0$) & 0.1822 & 0.0954 & 0.1980 & 0.1152 & 0.1760 & 0.0930 \\
\rowcolor{SoftBlue}
\textbf{RFT + Med-reward} & \textbf{0.1474} & \textbf{0.0865} & \textbf{0.1917} & \textbf{0.1055} & \textbf{0.1670} & \textbf{0.0892} \\
\bottomrule
\end{tabular}
}
}

\block{KEY FINDINGS}{
\begin{itemize}
    \item \textbf{In-domain gain:} on MultiMed, the proposed method reaches the best \textbf{WER = 0.1474} and \textbf{MWER = 0.0865}.
    \item \textbf{Cross-dataset robustness:} compared with SFT, MWER drops from \textbf{0.1323 to 0.1055} on Ankit (\textbf{20.3\% relative reduction}) and from \textbf{0.1149 to 0.0892} on Macabdul (\textbf{22.4\% relative reduction}).
    \item \textbf{Why it matters:} generic RFT alone is not enough; explicit medical-domain reward is the critical ingredient.
\end{itemize}
}

\block{CONCLUSION}{
\begin{itemize}
    \item Medical-aware RFT improves recognition of domain-critical terminology without sacrificing overall transcription quality.
    \item The method generalizes better than plain supervised fine-tuning on unseen medical datasets.
    \item \textbf{Future work:} leverage large-scale unlabeled medical audio and stronger medical NER / terminology grounding.
\end{itemize}

\textbf{Poster pitch tip:} prepare a 5--10 minute explanation focused on the key innovation, reward design, and cross-dataset results.
}

\end{columns}

\end{document}
